% ID: SEL-D0029-OA-E03
% ESTADO: cotejada. Segunda pasada visual completa y borrador compilado revisados.
% Fuente: Selectividad/matematicas_aplicadas/suplente septiembre/Examen-Matemáticas Aplicadas a las Ciencias Sociales.pdf
% SHA-256: 4b8884e307df4f5ddaf079e0f10be3d3f74d4793874eb03b355d34f4856579cb
% Página: 1; opción A; ejercicio original 3.

Se ha realizado un referéndum en el que se ha convocado a la ciudadanía a expresar con ``SÍ'' o con ``NO'' su opinión sobre cierta cuestión. En una determinada mesa electoral hay tres urnas que contienen las siguientes papeletas: la urna A tiene 200 papeletas con ``SÍ'' y 300 con ``NO'', la urna B, 500 ``SÍ'' y 400 ``NO'' y la urna C contiene 200 ``SÍ'' y 100 ``NO''.

Se elige una urna al azar y de ella se extrae aleatoriamente una papeleta.
\begin{enumerate}[label=\alph*)]
\item \textbf{(1.5 puntos)} Calcule la probabilidad de que sea un ``SÍ''.
\item \textbf{(1 punto)} Si la papeleta extraída es ``NO'', calcule la probabilidad de que haya sido extraída de la urna A.
\end{enumerate}
